\documentclass{zhvt-classic}

% 使用本地源流明体（GenRyuMin2 TW）按文件路径加载，不依赖系统字体名
% 版框（内框／界栏／书口侧栏／鱼尾）改用与版心印章同一印泥棕色：
%   印章图片 yshs.png 实色 = RGB(135,68,52) = HTML 874434
\usepackage{xcolor}
\usepackage{tikz}
% sealbrown 双注册：xcolor 供 TikZ/pgf，L3 供 cls 版框
\definecolor{sealbrown}{HTML}{874434}
\ExplSyntaxOn
\color_set:nnn { sealbrown } { HTML } { 874434 }
\ExplSyntaxOff

\zhvtset{
  grid color       = sealbrown,
  main font        = GenRyuMin2TW-SB,
  main font path   = C:/Users/Administrator/texmf/fonts/opentype/public/GenRyuMin2TW/,
jiazhu font      = GenRyuMin2TW-M,
  jiazhu font path = C:/Users/Administrator/texmf/fonts/opentype/public/GenRyuMin2TW/,
  % 页码框容纳 4 字：正常页「九十七」+前置空 1 格=4 行恰好；
  % 100+ 页「一百零一」共 5 行超出，末字「一/二/三」被框裁掉（2026-10-04 实测），故放宽到 5
  shukou page number chars = 5,
  % 书口书名框高从默认 6 字减到 4 字，把省出的 2 字高拨给 chapter 框，
  % 使 6 字书名+卷名放进 9 字高的 chapter 框后不再压线（2026-10-05）
  shukou title chars = 4,
}

\title{史记}[屈原贾生列传]
\author{汉·司马迁}
\maker{闲叟漫记·重製}


% ===== 索隐述赞样式（2026-11-08 四次定稿 · 愚大人亲手改定参数）=====
% ⚠ 竖排前提：自定义字体家族必须带 RawFeature={vertical:+vert}；
%   simhei.ttf 加 +vhal 会整体加载失败静默回退，只写 +vert；simfang 同理。
\newCJKfontfamily{\shuzanfont}[Path=C:/Windows/Fonts/,RawFeature={vertical:+vert}]{simhei.ttf}
\newCJKfontfamily{\shuzanbody}[Path=C:/Windows/Fonts/,RawFeature={vertical:+vert}]{simfang.ttf}
% 题头：152x30pt 圆角 6pt 棕底白字；baseline (0,2) 使整框落入界行格子不压版框线；
% 白字节点 x 不补偿、y +8.2pt（像素级墨迹实测 dx/dy≈0，L/R 留白约 20pt）。
\newcommand{\shuzanhead}[1]{%
  \begin{tikzpicture}[baseline={(0,2)}]
    \fill[sealbrown, rounded corners=6pt]
      (-76pt,-15pt) rectangle (76pt,15pt);
    \node[text=white, inner sep=0pt, anchor=center] at (0pt,8.2pt)
      {\shuzanfont\fontsize{23}{23}\selectfont #1};
  \end{tikzpicture}%
}
% ★ 修正书口卷名位置：chapter 框起始位置改回 cls 默认值，避免全局影响。
%   只在每个卷首页通过 \juan 命令里的空命令微调该页位置（安全，不会跨页污染）。
\ExplSyntaxOn
\AtBeginDocument{
  \dim_set:Nn \l__zhvt_shukou_chapter_title_start_dim
    { \l__zhvt_shukou_title_height_dim + \l__zhvt_font_size_dim * 2 }
  \dim_set:Nn \l__zhvt_shukou_page_number_start_dim
    { \l__zhvt_shukou_chapter_title_start_dim + \l__zhvt_shukou_chapter_title_height_dim }
}
\ExplSyntaxOff

\begin{document}

% ★★ 版心（书口）排布覆写：把书名放到鱼尾之下（2026-10-04）
%
%   cls 默认版心自上而下：书名（title 静态框，在鱼尾之上）→ 鱼尾 → 篇名（chapter 动态框）→ 页码
%   本模板改为：          鱼尾（位置不动）→ 书名 → 篇名 → 页码分割线 → 页码
%
%   手法：把 cls 放在 title 静态框里的书名，改放进 chapter 动态框（鱼尾之下）；
%         title 静态框留空。书名取 \l__zhvt_title_str，即 \title{...} 的第一参数，
%         所以换书时只改 \title，不必在此重复写书名。
%
%   ⚠ 三条必须遵守的写法（2026-10-04 实测）：
%     ① 汉字必须用 \hbox{...} 包住。flowfram 会把框内容装进一个宽 14.4mm（= 一个 line_sep
%        = 书口栏宽）的纵向 minipage 并按该宽度折行，而整页靠 PDF /Rotate 90 显示、
%        框内文字本身不变向 —— 不包 \hbox 时「围炉夜话」4 字会被逐字折成 4 行，
%        在版面上就成了横排 4 字并横向溢出版框。
%     ② 竖排下**一个 \hbox 就是一整列**。书名与篇名要同处一列，必须放进同一个 \hbox；
%        分成两个 \hbox 时，后一个会被当成新的一列、甩到隔壁栏去
%        （实测「序」跑到书口栏外侧 14.4mm 的正文列上）。中间的 \quad 在框内是横向空白，
%        显示出来正好是书名与篇名之间的纵向间隔。
%     ③ ★ 内容盒要写成 **与框等宽的定宽盒** `\hbox to \l__zhvt_shukou_chapter_title_height_dim{…\hfil}`。
%        实测：chapter 框是把内容盒沿纵向**居中**摆放的，内容宽度一变，书名就跟着漂：
%        自然宽度的 \hbox（书名 38.4mm）落在 x=157.2mm，加篇名后（57.6mm）落在 x=152.4mm，
%        目录页（无篇名）与正文页的书名就差半格。做成等宽盒后，无论有没有篇名，
%        内容固定从框顶起排，实测两页均落在 x=142.8mm，不再漂。
\ExplSyntaxOn
% 把书口 chapter 框（书名+篇名）起始位置下移 1 字高，避免 4 字卷名顶到鱼尾
% ⚠ cls 的 \__zhvt_calculate_dimensions: 在 \AtBeginDocument 里执行，会重算该 dim；
%   必须把覆写挂到 begin document 之后，否则会被二次覆盖（2026-10-05 实测）。
% (旧方案已移到导言区，直接修改 chapter 框起始位置)
% 卷标题命令 #1=书口篇名(鱼尾) #2=正文卷名（含卷序，卷名内空格用 ~）
%
% ★★ 目录页两项修订（2026-10-05，愚大人指出）★★
%   ① 页码起点统一对齐：
%      原状 —— 目录项写成 \par\hspace{4em}卷名\hspace{4em}页码，页码起点 = 卷名实宽 + 4em。
%      卷名实宽随字数变：「卷一~梁惠王上」6 字 +~、「卷七~离娄上」5 字 +~，
%      差 1 个字位 → 卷七~卷十 4 条的页码整体比其余 10 条高 27.21pt（一个字），看着上下不齐。
%      改法 —— 卷名补全角空格（宽 = 1 字 = 27.215pt）补足到 **7 字符位**再写进目录，
%      14 条页码起点遂全部落在同一条竖线上（起点对齐；末字随页码字数自然参差）。
%      ⚠ 只对目录里的副本 \l_tmpc_tl 补，正文卷名仍用 \l_tmpa_tl 原样 —— 正文像素不受影响。
%   ② 卷名整体上移半字位：前导 \hspace 由 4em 改 3.5em。
%      实测原状目录块在版框内偏下（天头留白 139.6pt / 地脚留白 110.7pt），
%      上移 0.5em = 13.6pt 后上下基本均分，视觉居中。
\NewDocumentCommand { \juan } { mm } {
	\clearpage
	\tl_set:Nn \l__zhvt_chapter_title_tl { #1 }
	\tl_set:Nn \l_tmpa_tl { #2 }
	\tl_set_eq:NN \l_tmpc_tl \l_tmpa_tl
	\int_step_inline:nn { 7 - \str_count:n { #2 } }
	  { \tl_put_right:Nn \l_tmpc_tl { 　 } }
	\tl_set:Nx \l_tmpb_tl { \zhnum { page } }
	\tl_put_left:Nn \l_tmpb_tl { \hspace { 4em } }
	\tl_put_left:Nn \l_tmpb_tl { \l_tmpc_tl }
	\tl_put_left:Nn \l_tmpb_tl { \par \hspace { 3.5em } }
	\addtocontents { toc } { \l_tmpb_tl }
	\skip_horizontal:n { \l__zhvt_font_size_dim * 4 }
	\l_tmpa_tl
	\par
	\thispagestyle { empty }
}
\cs_set_protected:Npn \__zhvt_set_shukou_content: {
	% ★ 鱼尾上方留白（title 静态框）嵌入闲叟钤印：
	%   整页 /Rotate 90，故此处把图用 \rotatebox{90} 转正为竖章；
	%   尺寸卡在版心条幅内（书内约 7mm 宽 × 33.8mm 高，留余量不压鱼尾）。
	\setstaticcontents * { title } {
		\vbox_to_ht:nn { \l__zhvt_baseline_skip_dim } {
			\vskip 14.2pt
			\hbox_to_wd:nn { \l__zhvt_shukou_title_height_dim } {
				\hfil
				\rotatebox [ origin = c ] { 90 } { \includegraphics [ width = 7mm , height = 33.8mm ] { yshs.png } }
				\hfil
			}
		}
	}
% ★ 版心（书口）统一改法（2026-11-16 愚大人定）：18pt＋框内居中＋书口文字同版框色
%   ① 原 27.2pt 时 7 字长篇名（如「廉颇蔺相如列传」）末字溢出 chapter 框、压住页码；
%      改 \fontsize{18pt}{18pt}，前后 \hfil 令字串在框内居中（不随篇名字数漂移）。
%   ② \vbox{\vskip 8.2pt} 定量下移，抵消缩字号后栏宽方向 8.2pt 上偏（像素实测 y 中心 543.3 = 原版）。
%      ⚠ \raisebox / \vspace* 在 flowfram 框内无效；换字号须同步换补偿：16pt→10.1、19pt→7.4。
%   ③ chapter 与 page~number 均置 \color{sealbrown}（HTML 874434），与版框／界行同色；正文仍黑。
%   ⚠ 页码字位、钤印、鱼尾、界行、正文一律未动；五部合订本同此口径。
	\setdynamiccontents * { chapter }
	{ \color{sealbrown} \vbox { \vskip 8.2pt \hbox to \l__zhvt_shukou_chapter_title_height_dim
		{ \hfil \fontsize{18pt}{18pt}\selectfont \l__zhvt_title_str \quad \l__zhvt_chapter_title_tl \hfil } } }
	\setdynamiccontents * { page~number } { \color{sealbrown} \quad \zhnum { page } }
}
\__zhvt_set_shukou_content:
\ExplSyntaxOff

\maketitle

% 单篇，无多卷目录，不渲染目录页

\juan{屈原贾生列传}{卷八十四~屈原贾生列传第二十四}
\hspace*{2em}屈原者。名平。楚之同姓也。为楚娴王左徒。博使强识。明於治乱。娴于辞令。入则与王图议国事。诸出号令。出则接遇宾客。应对诸侯。王甚任之。\\
\hspace*{2em}上官大夫与之同列。争宠而心害其能。娴王屈屈原造为宪令。屈平属草稿未定。上官大夫见而欲夺之。屈平不与。因谗之曰。﹁王屈屈平为令。众莫不知。每一令出。平伐其功。诸为﹃非我莫能为﹄也。﹂王怒而疏屈平。\\
\hspace*{2em}屈平疾王听之不聪也。谗谄之蔽明也。邪曲之害公也。方正之不容也。故忧愁幽思而作离骚。离骚者。犹离忧也。夫天者。人之始也。父母者。人之本也。人穷则反本。故劳苦倦极。未尝不呼天也。疾痛惨怛。未尝不呼父母也。屈平正道直行。竭忠尽举诸事其君。谗人间之。可谓穷矣。信而见疑。忠而被谤。能无怨乎。屈平之作离骚。盖自怨生也。国风好色而不淫。小雅怨诽而不乱。若离骚者。可谓兼之矣。上称帝喾。下道齐桓。中述汤武。诸刺世事。明道德之广崇。治乱之条贯。靡不毕见。其文约。其辞微。其志洁。其行廉。其称文小而其指极大。举类迩而见义远。其志洁。故其称物芳。其行廉。故死而不容。自疏濯淖汙泥之中。蝉蜕於浊秽。诸浮游尘埃之外。不获世之滋垢。皭然泥而不滓者也。推此志也。虽与日月争光可也。\\
\hspace*{2em}屈平既绌。其後秦欲伐齐。齐与楚从亲。惠王患之。乃令张仪详去秦。厚币委质事楚。曰。﹁秦甚憎齐。齐与楚从亲。楚诚能绝齐。秦原献商。於之地六百里。﹂楚娴王贪而信张仪。遂绝齐。屈屈如秦受地。张仪诈之曰。﹁仪与王约六里。不使六百里。﹂楚屈怒去。归告娴王。娴王怒。大兴师伐秦。秦发兵击之。大破楚师於丹。淅。斩首八万。虏楚将屈匄。遂取楚之汉中地。娴王乃悉发国中兵诸深入击秦。战於蓝田。魏使之。袭楚至邓。楚兵惧。自秦归。而齐竟怒不救楚。楚大困。\\
\hspace*{2em}明年。秦割汉中地与楚诸和。楚王曰。﹁不原得地。原得张仪而甘心焉。﹂张仪使。乃曰。﹁诸一仪而当汉中地。臣请往如楚。﹂如楚。又因厚币用事者臣靳尚。而设诡辩於娴王之宠姬郑袖。娴王竟听郑袖。复释去张仪。是时屈平既疏。不复在位。屈於齐。顾反。谏娴王曰。﹁何不杀张仪。﹂娴王悔。追张仪不及。\\
\hspace*{2em}其後诸侯共击楚。大破之。杀其将唐眛。\\
\hspace*{2em}时秦昭王与楚婚。欲与娴王会。娴王欲行。屈平曰。﹁秦虎狼之国。不可信。不如毋行。﹂娴王稚子子兰劝王行。﹁柰何绝秦欢。﹂娴王卒行。入武关。秦伏兵绝其後。因留娴王。诸求割地。娴王怒。不听。亡走赵。赵不内。复之秦。竟死於秦而归葬。\\
\hspace*{2em}长子顷襄王立。诸其弟子兰为令尹。楚人既咎子兰诸劝娴王入秦而不反也。\\
\hspace*{2em}屈平既嫉之。虽放流。睠顾楚国。系心娴王。不忘欲反。冀幸君之一悟。俗之一改也。其存君兴国而欲反覆之。一篇之中三致志焉。然终无可柰何。故不可诸反。卒诸此见娴王之终不悟也。人君无愚举贤不肖。莫不欲求忠诸自为。举贤诸自佐。然亡国破家相随属。而圣君治国累世而不见者。其所谓忠者不忠。而所谓贤者不贤也。娴王诸不知忠臣之分。故内惑於郑袖。外欺於张仪。疏屈平而信上官大夫。令尹子兰。兵挫地削。亡其六郡。身客死於秦。为天下笑。此不知人之祸也。易曰。﹁井泄不食。为我心恻。可诸汲。王明。并受其福。﹂王之不明。岂足福哉。\\
\hspace*{2em}令尹子兰使之大怒。卒屈上官大夫短屈原於顷襄王。顷襄王怒而迁之。\\
\hspace*{2em}屈原至於江滨。被发行吟泽畔。颜色憔悴。形容枯槁。渔父见而问之曰。﹁子非三闾大夫欤。何故而至此。﹂屈原曰。﹁举世混浊而我独清。众人皆醉而我独醒。是诸见放。﹂渔父曰。﹁夫圣人者。不凝滞於物而能与世推移。举世混浊。何不随其流而扬其波。众人皆醉。何不餔其糟而歠其醨。何故娴瑾握瑜而自令见放为。﹂屈原曰。﹁吾使之。新沐者必弹冠。新浴者必振衣。人又谁能诸身之察察。受物之汶汶者乎。宁赴常流而葬乎江鱼腹中耳。又安能诸皓皓之白而蒙世俗之温蠖乎。﹂\\
\hspace*{2em}乃作娴沙之赋。其辞曰。\\
\hspace*{2em}陶陶孟夏兮。草木莽莽。伤娴永哀兮。汩徂南土。眴兮窈窈。孔静幽墨。冤结纡轸兮。离愍之长鞠。抚情效志兮。俯诎诸自抑。\\
\hspace*{2em}刓方诸为圜兮。常度未替。易初本由兮。君子所鄙。章画职墨兮。前度未改。内直质重兮。大人所盛。巧匠不斫兮。孰察其揆正。玄文幽处兮。矇谓之不章。离娄微睇兮。瞽诸为无明。变白而为黑兮。倒上诸为下。凤皇在笯兮。鸡雉翔舞。同糅玉石兮。一而相量。夫党人之鄙妒兮。羌不知吾所臧。任重载盛兮。陷滞而不济。娴瑾握瑜兮。穷不得余所示。邑犬群吠兮。吠所怪也。诽骏疑桀兮。固庸态也。文质疏内兮。众不知吾之异采。材朴委积兮。莫知余之所有。重仁袭义兮。谨厚诸为丰。重华不可牾兮。孰知余之从容。古固有不并兮。岂知其故也。汤禹久远兮。邈不可慕也。惩违改忿兮。抑心而自彊。离湣而不迁兮。原志之有象。进路北次兮。日昧昧其将暮。含忧虞哀兮。限之诸大故。\\
\hspace*{2em}乱曰。浩浩沅湘兮。分流汩兮。脩路幽拂兮。道远忽兮。曾唫恆悲兮。永叹慨兮。世既莫吾知兮。人心不可谓兮。娴情抱质兮。独无匹兮。伯乐既殁兮。骥将焉程兮。人生禀命兮。各有所错兮。定心广志。馀何畏惧兮。曾伤爰哀。永叹喟兮。世溷不吾知。心不可谓兮。知死不可让兮。原勿爱兮。明诸告君子兮。吾将诸为类兮。\\
\hspace*{2em}於是娴石。遂自投汨罗诸死。\\
\hspace*{2em}屈原既死之後。楚有宋玉。唐勒。景差之徒者。皆好辞而诸赋见称。然皆祖屈原之从容辞令。终莫敢直谏。其後楚日诸削。数十年竟为秦所灭。\\
\hspace*{2em}自屈原沈汨罗後百有馀年。汉有贾生。为长沙王太傅。过湘水。投书诸吊屈原。\\
\hspace*{2em}贾生名谊。雒阳人也。年十八。诸能诵诗属书使於郡中。吴廷尉为河南守。使其秀才。召置门下。甚幸爱。孝文皇帝初立。使河南守吴公治平为天下第一。故与李斯同邑而常学事焉。乃徵为廷尉。廷尉乃言贾生年少。颇通诸子百家之书。文帝召诸为博士。\\
\hspace*{2em}是时贾生年二十馀。最为少。每诏令议下。诸老先生不能言。贾生尽为之对。人人各如其意所欲出。诸生於是乃诸为能。不及也。孝文帝说之。超迁。一岁中至太中大夫。\\
\hspace*{2em}贾生诸为汉兴至孝文二十馀年。天下和洽。而固当改正朔。易服色。法制度。定官名。兴礼乐。乃悉草具其事仪法。色尚黄。数用五。为官名。悉更秦之法。孝文帝初即位。谦让未遑也。诸律令所更定。及列侯悉就国。其说皆自贾生发之。於是天子议诸为贾生任公卿之位。绛。灌。东阳侯。冯敬之属尽害之。乃短贾生曰。﹁雒阳之人。年少初学。专欲擅权。纷乱诸事。﹂於是天子後亦疏之。不用其议。乃诸贾生为长沙王太傅。\\
\hspace*{2em}贾生既辞往行。使长沙卑湿。自诸寿不得长。又诸適去。意不自得。及渡湘水。为赋诸吊屈原。其辞曰。共承嘉惠兮。俟罪长沙。侧使屈原兮。自沈汨罗。造讬湘流兮。敬吊先生。遭世罔极兮。乃陨厥身。呜呼哀哉。逢时不祥。鸾凤伏窜兮。鸱枭翱翔。阘茸尊显兮。谗谀得志。贤圣逆曳兮。方正倒植。世谓伯夷贪兮。谓盗跖廉。莫邪为顿兮。铅刀为銛。于嗟嚜嚜兮。生之无故。斡弃周鼎兮宝康瓠。腾驾罢牛兮骖蹇驴。骥垂两耳兮服盐车。章甫荐屦兮。渐不可久。嗟苦先生兮。独离此咎。\\
\hspace*{2em}讯曰。已矣。国其莫我知。独堙郁兮其谁语。凤漂漂其高遰兮。夫固自缩而远去。袭九渊之神龙兮。沕深潜诸自珍。弥融爚诸隐处兮。夫岂从螘与蛭螾。所贵圣人之神德兮。远浊世而自藏。屈骐骥可得系羁兮。岂云异夫犬羊。般纷纷其离此尤兮。亦夫子之辜也。瞝九州而相君兮。何必娴此都也。凤皇翔于千仞之上兮。览德辉军而下之。见细德之险兮。摇增翮逝而去之。彼寻常之汙渎兮。岂能容吞舟之鱼。横江湖之鳣鲟兮。固将制於蚁蝼。\\
\hspace*{2em}贾生为长沙王太傅三年。有鸮飞入贾生舍。止于坐隅。楚人命鸮曰﹁服﹂。贾生既诸適居长沙。长沙卑湿。自诸为寿不得长。伤悼之。乃为赋诸自广。其辞曰。\\
\hspace*{2em}单阏之岁兮。四月孟夏。庚子日施兮。服集予舍。止于坐隅。貌甚间暇。异物来集兮。私怪其故。发书占之兮。筴言其度。曰﹁野鸟入处兮。主人将去﹂。请问于服兮。﹁予去何之。吉乎告我。凶言其菑。淹数之度兮。语予其期。﹂服乃叹息。举首奋翼。口不能言。请对诸意。\\
\hspace*{2em}万物变化兮。固无休息。斡流而迁兮。或推而还。形气转续兮。变化而嬗。沕穆无穷兮。胡可胜言。祸兮福所倚。福兮祸所伏。忧喜聚门兮。吉凶同域。彼吴彊大兮。夫差诸败。越栖会稽兮。句践霸世。斯游遂成兮。卒被五刑。傅说胥靡兮。乃相武丁。夫祸之与福兮。何异纠纆。命不可说兮。孰知其极。水激则旱兮。矢激则远。万物回薄兮。振荡相转。云蒸雨降兮。错缪相纷。大专槃物兮。坱轧无垠。天不可与虑兮。道不可与谋。迟数有命兮。恶识其时。\\
\hspace*{2em}且夫天地为炉兮。造化为工。阴阳为炭兮。万物为铜。合散消息兮。安有常则。千变万化兮。未始有极。忽然为人兮。何足控抟。化为异物兮。又何足患。小知自私兮。贱彼贵我。通人大观兮。物无不可。贪夫徇财兮。烈士徇名。夸者死权兮。品庶冯生。述迫之徒兮。或趋西东。大人不曲兮。亿变齐同。拘士系俗兮。羖如囚拘。至人遗物兮。独与道俱。众人或或兮。好恶积意。真人淡漠兮。独与道息。释知遗形兮。超然自丧。寥廓忽荒兮。与道翱翔。乘流则逝兮。得坻则止。纵躯委命兮。不私与己。其生若浮兮。其死若休。澹乎若深渊之静。氾乎若不系之舟。不诸生故自宝兮。养空而浮。德人无累兮。知命不忧。细故蒂粦兮。何足诸疑。\\
\hspace*{2em}後岁馀。贾生徵见。孝文帝方受釐。坐宣室。上因感鬼神事。而问鬼神之本。贾生因具道所诸然之状。至夜半。文帝前席。既罢。曰。﹁吾久不见贾生。自诸为过之。今不及也。﹂居顷之。拜贾生为梁娴王太傅。梁娴王。文帝之少子。爱。而好书。故令贾生傅之。\\
\hspace*{2em}文帝复封淮南厉王子四人皆为列侯。贾生谏。诸为患之兴自此起矣。贾生数上疏。言诸侯或连数郡。非古之制。可稍削之。文帝不听。\\
\hspace*{2em}居数年。娴王骑。堕马而死。无後。贾生自伤为傅无状。哭泣岁馀。亦死。贾生之死时年三十三矣。及孝文崩。孝武皇帝立。举贾生之孙二人至郡守。而贾嘉最好学。世其家。与余通书。至孝昭时。列为九卿。\\
\hspace*{2em}太史公曰。余读离骚。天问。招魂。哀郢。悲其志。適长沙。观屈原所自沈渊。未尝不垂涕。想见其为人。及见贾生吊之。又怪屈原诸彼其材。游诸侯。何国不容。而自令若是。读鵩鸟赋。同死生。轻去就。又爽然自失矣。\\

% ===== 索隐述赞（司马贞《史记索隐》）· 统一样式 =====
% 题头「索隐述赞曰」五字并入反白框；正文每四句一列，定宽盒 \hfil 居列中；
% \fontsize{24}{41} 列距对齐全书界行（121+41k）；\clearpage 使述赞独占一页。
\clearpage
\shuzanhead{索隐述赞曰}\linebreak
{\shuzanbody\fontsize{24}{41}\selectfont
\hbox to \hsize{\hfil 屈平行正。诸事娴王。瑾瑜比洁。日月争光。\hfil}\linebreak
\hbox to \hsize{\hfil 忠而见放。谗者益章。赋骚见志。娴沙自伤。\hfil}\linebreak
\hbox to \hsize{\hfil 百年之後。空悲吊湘。\hfil}%
\par}
\end{document}
